\documentclass[7pt]{article}
%\usetheme{Warsaw}

%\usepackage{beamerinnerthemecircles, beamerouterthemeshadow}
%\usepackage{scrextend}

%\changefontsizes{7pt}


% Standard packages
\usepackage{graphicx}
\usepackage{graphics}

\usepackage{color}
\usepackage[OT4]{fontenc}
\usepackage[utf8]{inputenc}
%\usepackage[latin1]{inputenc}
\usepackage{times}\usepackage{amsmath}
\usepackage[T1]{fontenc}

%\mode<presentation> {
%    \usetheme{Singapore} %Bergen, Berkely, Berlin, Boadilla, CambridgeUS, Darmstadt,
%                          %Frankfurt, Goettingen, Singapore, Warsaw
%    \usecolortheme{default} %beetle, seahorse, wolverine, dolphin, beaver
%}


\title{Wave equation for a 2D membrane (FEM version)}

\definecolor{whi}{rgb}{0.95,0.95,0.95}
% Setup TikZ

%\usepackage{tikz}
%\usetikzlibrary{arrows}
%\tikzstyle{block}=[draw opacity=0.7,line width=1.4cm]


% Author, Title, etc.




% The main document

\begin{document}
\maketitle




\section{the problem}


We solve the advection-diffusion equation
\begin{equation} \frac{\partial u}{\partial t}=-v_x\frac{\partial u}{\partial x} - v_y\frac{\partial u}{\partial y} +D\nabla^2 u, \end{equation}
with the velocity field $(v_x,v_y)$ and the diffusion constant independent of the position.
We use a square computational box $(x,y)\in(0,L)\times(0,L)$ and the finite element method and introduce a periodic boundary conditions
\begin{equation} u(x+mL,y+nL)=u(x,y) \end{equation} for integer $m$ and $n$. In FEM application of the periodic boundary conditions is easy:
the nodes from the edges of the box should be attributed the same global number, e.g. the node at the upper edge $(x,L)$ gets the same global number as the node
in the position $(x,0)$ etc. 

In our previous assignments we had evaluated the global matrix of $-\nabla^2$ (the stiffness matrix {\bf S}).
We need to evaluate the matrix of the advection operator. For a constant velocity the local matrix
is the same for all the elements. The local matrix for the first derivative in $x$ direction is evaluated as  
$$(\frac{\partial}{\partial x})_{ji}=\frac{a}{2}\sum_{l}\sum_{n}  w_l w_n  g_j(p_l,p_n)\frac{ g_i(p_l+\Delta,p_n)-g_i(p_l-\Delta,p_n)}{2\Delta }.$$
We evaluate similarily $(\frac{\partial}{\partial y})$ matrix and the assembled global matrix of the advection operator is denoted by ${\bf C}$,
and the overlap matrix is denoted by ${\bf O}$

The Crank-Nicolson step is calculated as the solution of the linear system of equations
$$\left({\bf O} -\frac{\Delta t}{2}({\bf C-S})\right){\bf c}(t+\Delta t)=\left({\bf O} +\frac{\Delta t}{2}({\bf C-S})\right){\bf c}(t).$$


Lets set $\Delta t=0.02$, $L=5$ and the initial condition:   $u(x,y)=\exp(-2(x-L/2)^2-2(y-L/2)^2)$.




%where ${\bf H}(t)={\bf H}_0+{\bf V} \sin(\omega t)$.


\section{Deliverables}
\begin{itemize}
\item pure advection $D=0$, $v_x=v_y=1$. Solve the equations till the center of the density returns to its initial position second time. Is its shape constant?
\item pure diffusion $D=0.1$, $v_x=v_y=0$.  Plot the minimal and maximal values of $u$ at the nodes as a function of time.
\item advection diffusion $D=0.1$, $v_x=v_y=1$. Is the velocity of the center of the packet equal to the velocity field vector?
\end{itemize}
\end{document}




